%==============================================================================
%  Appendix: Fundamentals of Hamiltonian Monte Carlo
%  Companion to "Multilevel Models for Airbnb Ratings"
%
%  Standalone: pdflatex appendix-hmc.tex
%  Or drop into report.tex with \input{appendix-hmc-body} after \appendix.
%==============================================================================
\documentclass[10pt,a4paper]{article}

\usepackage[T1]{fontenc}
\usepackage[utf8]{inputenc}
\usepackage[top=0.65in,bottom=0.7in,left=0.85in,right=0.85in]{geometry}
\usepackage{amsmath,amssymb,amsthm}
\usepackage{enumitem}
\usepackage[compact]{titlesec}

\titlespacing*{\section}{0pt}{6pt}{3pt}

\setlist{nosep,leftmargin=*}
\setlength{\parskip}{0pt}

% the document is mostly displayed mathematics; the default skips around
% displays are what push it past two pages
\setlength{\abovedisplayskip}{4pt}
\setlength{\belowdisplayskip}{4pt}
\setlength{\abovedisplayshortskip}{2pt}
\setlength{\belowdisplayshortskip}{2pt}

\theoremstyle{definition}
\newtheorem*{defn}{Definition}
\newtheorem*{thm}{Theorem}
\newtheorem*{prop}{Proposition}

% compact framed block for the algorithm
\newenvironment{algobox}[1]{%
  \par\vspace{3pt}\begin{center}\begin{minipage}{0.97\textwidth}%
  \rule{\linewidth}{0.8pt}\vspace{1pt}\par\textbf{#1}\par\vspace{1pt}\small
}{%
  \par\vspace{1pt}\rule{\linewidth}{0.8pt}\end{minipage}\end{center}\vspace{3pt}\par
}

\title{\vspace{-2em}\textbf{Appendix: Fundamentals of Hamiltonian Monte Carlo
       for Bayesian Posterior Sampling}}
\author{}
\date{}

\begin{document}
\maketitle
\vspace{-3em}

%==============================================================================
\section{From Markov Chains to Metropolis--Hastings}
%==============================================================================

We wish to draw from a target $\pi(x) = f(x)/Z$ on $\mathbb{R}^D$, where $f$ can
be evaluated pointwise but the normaliser $Z = \int f(x)\,dx$ cannot. MCMC gives
up on sampling $\pi$ directly and instead builds a Markov chain that visits
states in proportion to $\pi$.

\begin{defn}[Markov chain and invariance]
A Markov chain on $\mathbb{R}^D$ is specified by a transition kernel
$T(x \to x')$, the density of moving to $x'$ from $x$. A distribution $\pi$ is
\emph{invariant} (stationary) for $T$ if
$\int \pi(x)\, T(x \to x')\, dx = \pi(x')$. The design problem is to construct a
$T$ whose invariant distribution is the $\pi$ we want.
\end{defn}

\paragraph{Metropolis--Hastings.} From the current state $x$, draw a proposal
$x^{*} \sim q(x^{*} \mid x)$ and accept it with probability
\[
  \alpha(x \to x^{*})
  = \min\!\left(1,\; \frac{f(x^{*})\, q(x \mid x^{*})}
                          {f(x)\, q(x^{*} \mid x)}\right),
\]
otherwise remain at $x$. The target enters only through the \emph{ratio}
$f(x^{*})/f(x)$, so the intractable $Z$ cancels. This is the single fact that
makes Bayesian computation possible at all.

\begin{defn}[Detailed balance]
$T$ satisfies detailed balance with respect to $\pi$ if
$\pi(x)\, T(x \to x') = \pi(x')\, T(x' \to x)$ for all $x, x'$.
\end{defn}

\begin{prop}
Detailed balance implies invariance: integrating both sides over $x$ gives
$\int \pi(x) T(x \to x')\,dx = \pi(x') \int T(x' \to x)\,dx = \pi(x')$, since
$T(x' \to \cdot)$ is a density. Metropolis--Hastings is constructed to satisfy
detailed balance, which is why it is correct for \emph{any} proposal $q$.
\end{prop}

\paragraph{The curse of dimensionality.} Correct is not efficient. In high
dimension the posterior mass does not sit near the mode but concentrates on a
thin shell --- the \emph{typical set} --- at radius $O(\sqrt{D})$, since the
density falls off while the volume element grows like $r^{D-1}$. An isotropic
random-walk proposal $\mathcal{N}(x, \epsilon^{2}I)$ almost always steps off
that shell into negligible density, so acceptance decays exponentially in $D$
unless $\epsilon \propto D^{-1/2}$. Shrinking $\epsilon$ makes the chain
diffuse: the iterations needed per effectively independent draw grow like
$O(D)$. The sampler can only shuffle along a shell it should be traversing.

%==============================================================================
\section{Hamiltonian Monte Carlo}
%==============================================================================

HMC removes the random walk by giving the chain \emph{momentum}, so that a
single proposal can traverse the typical set rather than jitter within it.

\paragraph{Phase space.} Augment the position $\boldsymbol{\theta} \in
\mathbb{R}^{D}$ with an auxiliary momentum $\mathbf{p} \sim
\mathcal{N}(\mathbf{0}, \mathbf{M})$, giving a $2D$-dimensional phase space.
Define the \emph{Hamiltonian}
\[
  H(\boldsymbol{\theta}, \mathbf{p}) = U(\boldsymbol{\theta}) + K(\mathbf{p}),
  \qquad
  U(\boldsymbol{\theta}) = -\ln f(\boldsymbol{\theta}),
  \qquad
  K(\mathbf{p}) = \tfrac{1}{2}\mathbf{p}^{\!\top}\mathbf{M}^{-1}\mathbf{p}.
\]
The joint density $\pi(\boldsymbol{\theta}, \mathbf{p}) \propto
\exp\{-H(\boldsymbol{\theta}, \mathbf{p})\}$ factorises, so its
$\boldsymbol{\theta}$-marginal is exactly the target. Momentum is scaffolding:
it is resampled each iteration and discarded.

\paragraph{Intuition.} $U$ is a landscape whose valleys are regions of high
posterior density, and the state is a frictionless puck on it. Kicked in a
random direction, the puck climbs slopes and turns back under its own momentum,
sweeping along a contour of constant total energy --- covering long distances
\emph{within} the typical set instead of jittering across it.

\begin{defn}[Hamilton's equations]
\[
  \frac{d\boldsymbol{\theta}}{dt} = \frac{\partial H}{\partial \mathbf{p}}
    = \mathbf{M}^{-1}\mathbf{p},
  \qquad
  \frac{d\mathbf{p}}{dt} = -\frac{\partial H}{\partial \boldsymbol{\theta}}
    = -\nabla U(\boldsymbol{\theta}).
\]
Write $\Phi_{t}$ for the flow map carrying $(\boldsymbol{\theta}(0),
\mathbf{p}(0))$ to $(\boldsymbol{\theta}(t), \mathbf{p}(t))$.
\end{defn}

\begin{thm}[Conservation of energy]
Along any solution, $\frac{dH}{dt} = 0$; hence
$H(\boldsymbol{\theta}(t), \mathbf{p}(t)) =
 H(\boldsymbol{\theta}(0), \mathbf{p}(0))$ for all $t$.
\end{thm}

\begin{thm}[Liouville's theorem]
The vector field of Hamilton's equations is divergence-free, so $\Phi_{t}$
preserves phase-space volume: $\det J_{\Phi_{t}} = 1$ for all $t$.
\end{thm}

These two facts are what a Metropolis correction needs. A proposal generated by
running $\Phi_{L}$ then negating the momentum is a deterministic involution,
hence reversible, with acceptance ratio
\[
  \alpha = \min\!\Big(1,\;
    \exp\{-H(\boldsymbol{\theta}^{*}, \mathbf{p}^{*})
          + H(\boldsymbol{\theta}_{0}, \mathbf{p}_{0})\}
    \cdot \big|\det J_{\Phi_{L}}\big| \Big).
\]
Liouville's theorem sets the Jacobian to $1$; conservation of energy sets the
exponent to $0$. So an \emph{exact} Hamiltonian trajectory is accepted with
probability $1$, no matter how far it travels.

%==============================================================================
\section{Bayesian Inference and the Complete Algorithm}
%==============================================================================

\paragraph{The posterior challenge.} For parameters $\boldsymbol{\theta}$ and
data $\mathcal{D}$,
\[
  p(\boldsymbol{\theta} \mid \mathcal{D})
  = \frac{p(\mathcal{D} \mid \boldsymbol{\theta})\, p(\boldsymbol{\theta})}
         {\int p(\mathcal{D} \mid \boldsymbol{\theta})\,
                p(\boldsymbol{\theta})\, d\boldsymbol{\theta}}
  = \frac{f(\boldsymbol{\theta})}{Z}.
\]
The evidence $Z = p(\mathcal{D})$ is a $D$-dimensional integral: grid quadrature
costs $O(k^{D})$, and conjugate updates do not exist for hierarchical models
with non-conjugate likelihoods. But HMC needs only
\[
  U(\boldsymbol{\theta}) = -\ln p(\mathcal{D} \mid \boldsymbol{\theta})
                           - \ln p(\boldsymbol{\theta}),
\]
and its gradient, both computable without $Z$.

\paragraph{Leapfrog integration.} Hamilton's equations rarely have closed-form
solutions, so the flow is approximated in steps of size $\epsilon$ by the
\emph{leapfrog} (Störmer--Verlet) scheme --- half a momentum kick, a full
position drift, and a second half kick:
\begin{align*}
  \mathbf{p}(t + \tfrac{\epsilon}{2})
    &= \mathbf{p}(t) - \tfrac{\epsilon}{2}\nabla U(\boldsymbol{\theta}(t)),\\
  \boldsymbol{\theta}(t + \epsilon)
    &= \boldsymbol{\theta}(t)
       + \epsilon\, \mathbf{M}^{-1}\mathbf{p}(t + \tfrac{\epsilon}{2}),\\
  \mathbf{p}(t + \epsilon)
    &= \mathbf{p}(t + \tfrac{\epsilon}{2})
       - \tfrac{\epsilon}{2}\nabla U(\boldsymbol{\theta}(t + \epsilon)).
\end{align*}

\begin{thm}[Symplecticity and the shadow Hamiltonian]
Each leapfrog step is a composition of shear maps, so it preserves phase volume
\emph{exactly} ($\det J = 1$) and is time-reversible. Being symplectic, it
conserves exactly a nearby \emph{shadow Hamiltonian}
$\tilde{H} = H + \mathcal{O}(\epsilon^{2})$, so energy error stays bounded over
arbitrarily long trajectories instead of accumulating. Forward Euler, by
contrast, has $\det J \neq 1$ and no shadow Hamiltonian: its energy drifts
monotonically and trajectories spiral away.
\end{thm}

Because volume preservation and reversibility hold \emph{exactly} for leapfrog,
the Metropolis correction below restores exactness for any $\epsilon$; the step
size affects only efficiency, through the acceptance rate.

\begin{algobox}{Hamiltonian Monte Carlo}
\textbf{Input:} initial state $\boldsymbol{\theta}^{(0)}$, trajectory length
$L$, step size $\epsilon$, mass matrix $\mathbf{M}$.

\smallskip
\textbf{For} $s = 1, \dots, S$:
\begin{enumerate}
  \item Sample fresh momentum
        $\mathbf{p}^{(s)} \sim \mathcal{N}(\mathbf{0}, \mathbf{M})$.
  \item Set $(\boldsymbol{\theta}_{0}, \mathbf{p}_{0})
        = (\boldsymbol{\theta}^{(s-1)}, \mathbf{p}^{(s)})$.
  \item Run $L$ leapfrog steps of size $\epsilon$ using
        $\nabla U(\boldsymbol{\theta})$ to obtain
        $(\boldsymbol{\theta}^{*}, \mathbf{p}^{*})$.
  \item Compute
        $\alpha = \min\!\big(1, \exp\{-H(\boldsymbol{\theta}^{*}, \mathbf{p}^{*})
        + H(\boldsymbol{\theta}_{0}, \mathbf{p}_{0})\}\big)$.
  \item With probability $\alpha$ set
        $\boldsymbol{\theta}^{(s)} = \boldsymbol{\theta}^{*}$;
        otherwise $\boldsymbol{\theta}^{(s)} = \boldsymbol{\theta}^{(s-1)}$.
\end{enumerate}
\textbf{Output:} draws $\boldsymbol{\theta}^{(1)}, \dots,
\boldsymbol{\theta}^{(S)}$.
\end{algobox}

\paragraph{From draws to summaries.} The output is an $S \times D$ matrix whose
rows are dependent but marginally distributed as the posterior.

\begin{thm}[Ergodic theorem for MCMC]
If the chain is irreducible and aperiodic with invariant distribution
$p(\boldsymbol{\theta} \mid \mathcal{D})$, then for any integrable $g$,
\[
  \frac{1}{S}\sum_{s=1}^{S} g\big(\boldsymbol{\theta}^{(s)}\big)
  \;\xrightarrow[S \to \infty]{\text{a.s.}}\;
  \mathbb{E}_{p(\boldsymbol{\theta} \mid \mathcal{D})}
    \big[g(\boldsymbol{\theta})\big].
\]
\end{thm}

This licenses every number in the main report. For the multilevel model
$\alpha_{j} \sim \mathrm{Normal}(\bar\alpha, \sigma)$, a listing's posterior
mean intercept is the column average of that $\alpha_{j}$, an $89\%$ credible
interval is the corresponding pair of column quantiles, and a derived quantity
--- $\mathrm{logit}^{-1}(\alpha_{j})$, or a prediction for an unseen listing
$\mathrm{logit}^{-1}(\bar\alpha + z\sigma)$ --- is estimated by transforming
row-wise and \emph{then} averaging. The order matters: averaging first
estimates a different quantity, since $\mathrm{logit}^{-1}$ is non-linear.

Two caveats follow. The guarantee is asymptotic, so $n_{\text{eff}}$ must be
reported alongside it. And a \emph{divergent transition} signals that leapfrog
has broken down: where the posterior curves more sharply than $\epsilon$ can
resolve, $\tilde{H}$ stops tracking $H$ and the region is under-explored, so
divergences indicate bias rather than mere inefficiency --- which is why the
main report reparameterises instead of ignoring them.

\end{document}
